\section{The ceiling}
\label{sec:vsceiling}

A fifth of the fair strike accrues where nothing is quoted.  Before
deciding what a model \emph{should} draw there, establish what any
model \emph{may} draw there.  This section proves the one bound that
holds with no data at all: a universal ceiling on how steeply total
variance can grow with strike.

Recall the objects.  The wing slopes of a smile are its asymptotic
growth rates in total variance,
\[
\beta_R=\limsup_{k\to+\infty}\frac{w(k)}{k},
\qquad
\beta_L=\limsup_{k\to-\infty}\frac{w(k)}{|k|},
\]
the same $\beta_L,\beta_R$ carried since \cref{sec:tails} (a limsup is
the largest value a ratio keeps returning to; for every smile in this
book the ratios simply converge).  No admissible smile --- no smile
that prices a positive law with mean one --- can have either slope
above $2$.

The proof rests on one physical fact and one computation.  The fact:
a call at an absurdly high strike must be worth almost nothing.
Formally, $c(k)=\E[(Y-e^{k})^{+}]$, the payoff is dominated by $Y$
itself (which has finite mean) and falls to zero pointwise as
$k\to\infty$, so dominated convergence gives
\begin{equation}
c(k)\longrightarrow0
\qquad(k\to+\infty).
\label{eq:vscalldies}
\end{equation}
The computation: what the Black function does when total variance
grows along a ray.  Fix a slope $\beta>0$ and evaluate
$\Black(k,\beta k)$ for large $k$.  The Black arguments (from
\cref{eq:black}) are $d_{\pm}=-k/\sqrt{w}\pm\sqrt{w}/2$; substituting
$w=\beta k$,
\begin{equation}
d_{+}=\sqrt{k}\left(\frac{\sqrt\beta}{2}-\frac{1}{\sqrt\beta}\right),
\qquad
d_{-}=-\sqrt{k}\left(\frac{1}{\sqrt\beta}+\frac{\sqrt\beta}{2}\right).
\label{eq:vsray}
\end{equation}
Everything hinges on the sign of the parenthesis in $d_{+}$:
multiplying $\sqrt\beta/2>1/\sqrt\beta$ through by $\sqrt\beta$ shows
it is positive exactly when $\beta>2$.  So as $k\to\infty$, $d_{+}$
runs to $+\infty$ for $\beta>2$, to $-\infty$ for $\beta<2$, and sits
at exactly $0$ for $\beta=2$.

The second Black term dies in every case.  The tool is the normal
tail bound $\PhiN(d)\le\phin(d)/|d|$ for $d<0$, with $\PhiN,\phin$
the standard normal distribution function and density.  Its one-line
proof: for $t\le d<0$ the ratio $t/d$ is at least $1$, so
\[
\PhiN(d)=\int_{-\infty}^{d}\phin(t)\,\dd t
\le\int_{-\infty}^{d}\frac{t}{d}\,\phin(t)\,\dd t
=\frac{\phin(d)}{|d|},
\]
the last step because $t\phin(t)$ integrates to $-\phin(d)$.
Applying the bound,
\[
\log\big(e^{k}\PhiN(d_{-})\big)
\le k-\frac{d_{-}^{2}}{2}-\log\big(|d_{-}|\sqrt{2\pi}\big),
\qquad
\frac{d_{-}^{2}}{2}
=\frac{k}{2}\left(\frac{1}{\beta}+1+\frac{\beta}{4}\right),
\]
and the coefficient of $k$ collects to
\begin{equation}
1-\frac12\left(\frac1\beta+1+\frac\beta4\right)
=\frac{4\beta-4-\beta^{2}}{8\beta}
=-\frac{(\beta-2)^{2}}{8\beta}\;\le\;0 ,
\label{eq:vsdecay}
\end{equation}
strictly negative except at $\beta=2$, where the $\log|d_{-}|$ term
still drives the bound to $-\infty$.  Hence
$e^{k}\PhiN(d_{-})\to0$ along every ray, and
\begin{equation}
\Black(k,\beta k)=\PhiN(d_{+})-e^{k}\PhiN(d_{-})
\;\longrightarrow\;
\begin{cases}
0, & \beta<2,\\[2pt]
\tfrac12, & \beta=2,\\[2pt]
1, & \beta>2 .
\end{cases}
\label{eq:vsraylimits}
\end{equation}

\begin{proposition}[The model-free ceiling]
\label{prop:vsceiling}
For any admissible smile, $\beta_L\le2$ and $\beta_R\le2$.
\end{proposition}

\begin{proof}
Right wing.  Suppose $\beta_R>2$ and pick $\beta$ strictly between $2$
and $\beta_R$.  By the definition of the limsup there are arbitrarily
large strikes with $w(k)\ge\beta k$.  The Black function increases in
its variance argument, so at those strikes
$c(k)=\Black(k,w(k))\ge\Black(k,\beta k)$, and by
\eqref{eq:vsraylimits} the right side approaches $1$.  A smile that
steep prices calls that never become worthless, contradicting
\eqref{eq:vscalldies}.

Left wing.  The mirror argument runs through the put, measured per
unit of strike: $P(k)/e^{k}=\E[(1-Y e^{-k})^{+}]$, which by dominated
convergence tends to $\Prob(Y=0)$ as $k\to-\infty$ --- and
$\Prob(Y=0)=1$ is impossible for a law with $\E[Y]=1$.  It remains to
show a left ray of slope $\beta>2$ forces that limit to be $1$.
Parity turns the Black call into the Black put per unit of strike,
\[
\frac{P(k)}{e^{k}}
=\frac{\Black(k,w)-1+e^{k}}{e^{k}}
=\PhiN(-d_{-})-e^{-k}\,\PhiN(-d_{+}),
\]
and along $w=\beta|k|$ (so $k=-|k|$) the arguments
\eqref{eq:vsray} give
\[
-d_{-}=\sqrt{|k|}\left(\frac{\sqrt\beta}{2}-\frac{1}{\sqrt\beta}\right),
\qquad
-d_{+}=-\sqrt{|k|}\left(\frac{1}{\sqrt\beta}+\frac{\sqrt\beta}{2}\right):
\]
exactly the two expressions of the right-wing computation with the
roles of the arguments exchanged.  So for $\beta>2$ the first term
$\PhiN(-d_{-})\to1$, and the second, $e^{|k|}\PhiN(-d_{+})$, dies by
the same exponent $-(\beta-2)^{2}/(8\beta)$ as in
\eqref{eq:vsdecay}.  Hence $P(k)/e^{k}\to1$, the contradiction.
\end{proof}

\begin{figure}[htbp]
\centering
\paperfig{0.9\textwidth}{fig_vs_ceiling.pdf}
\caption{The ceiling, computed.  The normalized Black call along
total-variance rays $w=\beta k$, against log-strike on a log axis.
Below the ceiling the call dies (at $k=400$ the $\beta=1.6$ ray reads
\MacVsRayBelowLim); above it the call saturates at forward value
(\MacVsRayAboveLim\ for $\beta=2.4$); on the boundary ray it creeps up
to one half (\MacVsRayAtLim\ at the right edge).  Only the dying
curves are smiles of any law.}
\label{fig:vsceiling}
\end{figure}

\Cref{fig:vsceiling} draws the computation.  Each curve is
$\Black(k,\beta k)$ for one slope, out to $k=400$ on a logarithmic
strike axis --- far past anything tradeable, because the convergence
is slow and showing it takes that range.  The three regimes of
\eqref{eq:vsraylimits} separate cleanly: the $\beta<2$ curves sink to
zero (the $\beta=1.6$ ray reads \MacVsRayBelowLim\ at the right
edge), the $\beta>2$ curves climb toward one (\MacVsRayAboveLim\ for
$\beta=2.4$), and the dashed boundary ray $\beta=2$ creeps up to one
half --- there $\PhiN(d_{+})$ equals $\tfrac12$ exactly at every
strike, and what remains is the second term dying at its slow
boundary rate, leaving \MacVsRayAtLim\ at the right edge.  Notice
that the boundary ray also violates
\eqref{eq:vscalldies}: a call permanently worth half the forward is no
better than one worth the whole forward.  The ceiling is therefore a
\emph{supremum, not a member}: an admissible smile can approach slope
$2$ from below, but no admissible smile lies on the ray itself.
Chapter~3 met the same phenomenon inside the density factor: at
$\beta=2$ the leading tail term of $g_{\rm D}$ vanishes, and whether
the far density is positive is decided by the next, smaller term ---
one the slope condition never examines --- which is why the working
cap of \cref{eq:leecap} sits strictly inside the bound.
\Cref{sec:vsenvelope} will meet the same phenomenon a third time, as
the top edge of the envelope.

The ceiling says how steep a wing \emph{cannot} be.  Lee's moment
formula --- stated in Chapter~2 as \cref{eq:lee} --- says how steep it
\emph{is}: $\beta=\Psi(r^{*})$, with
$\Psi(r)=2-4(\sqrt{r^{2}+r}-r)$, where $r^{*}$ is the tail's critical
moment order and $\Psi$ decreases from $\Psi(0)=2$.  The two results
are faces of one mechanism, and the computation above quietly
contained the link.  For $\beta<2$ the surviving Black term is
$\PhiN(d_{+})$ with $d_{+}\to-\infty$, and the same tail bound gives
its decay exponent:
\[
\frac{d_{+}^{2}}{2}
=\frac{k}{2}\left(\frac{1}{\sqrt\beta}-\frac{\sqrt\beta}{2}\right)^{2}
=\frac{k\,(2-\beta)^{2}}{8\beta}
\;\Longrightarrow\;
\PhiN(d_{+})\approx e^{-k\,(2-\beta)^{2}/(8\beta)}=K^{-r},
\quad
r=\frac{(2-\beta)^{2}}{8\beta},
\]
up to algebraic factors --- and this $r$ is exactly the moment budget
$\Psi^{-1}(\beta)$ that Chapter~3 computed in \cref{eq:momentmap}.
The bridge between call decay and moments is elementary: whenever
$Y\ge K$ one has $(Y-K)^{+}\le Y\le Y\,(Y/K)^{r}$, so a finite moment
caps the call, $\E[(Y-K)^{+}]\le\E[Y^{1+r}]\,K^{-r}$ --- more finite
moments force faster call decay, and (this is Lee's harder half) the
converse holds too.  A wing of slope
$\beta$ prices calls that decay with power $r$; call decay and finite
moments carry the same information; so slope pins moments and moments
pin slope.  Steeper wings, fatter tails, fewer moments --- one dial.

What mathematics gives, then, is a cone: every admissible wing lies
between slope $0$ and slope $2$, and its slope names its tail's moment
order.  What mathematics does not give is a selection.  The next
section sharpens the cone into a finite-strike envelope anchored at
the last quote --- and shows that it, too, bounds without selecting.
